% NSF Synergistic Activities (PAPPG 24-1, II.D.2.h(iv)) -- DRAFT
%
% Required for each senior person (here, only the PI) and uploaded in Research.gov as a
% separate PDF; unlike the biosketch, it is not made in SciENcv. Limit: 1 page, with up to
% five distinct examples that demonstrate the broader impact of the individual's professional
% and scholarly activities (creation, integration, and transfer of knowledge). PAPPG examples
% include innovations in teaching and training, development of research tools, broadening
% participation of groups underrepresented in STEM, and service to the scientific community
% outside the individual's own organization.
%
% Items are drawn from the CV, the Project Description, and the 2025 MSCS award application
% (the source of the geometry advisor panel). No public notes exist from these seminars.

\documentclass[11pt]{article}

\usepackage[margin=1in]{geometry}
\usepackage{amsmath, amssymb}
\usepackage{microtype}
\usepackage{enumitem}
\usepackage[compact]{titlesec}

\titleformat{\section}{\large\bfseries}{\thesection.}{0.5em}{}
\titlespacing*{\section}{0pt}{1.2ex plus .4ex minus .2ex}{0.5ex plus .2ex}
\setlist{nosep, leftmargin=*}

% Block paragraphs, as in the Project Description.
\setlength{\parindent}{0pt}
\setlength{\parskip}{0.3\baselineskip plus 2pt minus 1pt}

\pagestyle{empty}

\begin{document}

\begin{center}
  {\Large\bfseries Synergistic Activities}\\[5pt]
  Edward A. Varvak
\end{center}

\begin{enumerate}[label=\arabic*., itemsep=0.5\baselineskip]

\item \textbf{Seminar organization.} Co-organizer of the Graduate Algebraic Geometry Seminar
at the University of Illinois Chicago (UIC), Fall 2025--Spring 2026, a student-run seminar in
which graduate students give talks on algebraic geometry.

\item \textbf{Learning seminars.} Organizer of two learning seminars for graduate students at
UIC: a reading group on the Weil conjectures (Fall 2025) and a seminar on Lie groups (Fall
2024). Speaker at the UIC Graduate Student Colloquium (February 2025), presenting research
to a general audience of graduate students.

\item \textbf{Mathematics outreach.} Judge at QED, Chicago's Youth Math Symposium, run by the
Math Circles of Chicago (December 2025), at which students in grades 5--12 present original
projects in mathematics and computer science.

\item \textbf{Teaching.} Graduate student lecturer for Calculus~II at the University of
Chicago (Fall 2026). Teaching assistant for five semesters at UIC, a Hispanic-Serving
Institution where many students are the first in their families to attend college, in
Finite Mathematics for Business (four semesters) and Calculus~II (one semester).

\item \textbf{Mentoring graduate students.} Member of the geometry advisor panel of the UIC
Department of Mathematics, Statistics, and Computer Science, and an informal mentor to
younger graduate students navigating graduate school.

\end{enumerate}

\end{document}
